\lysh{22192_SOLUTION}{1}
\textnormal{ΛΥΣΗ}

\textnormal{α) Η έλλειψη με εξίσωση την (1) έχει } $\alpha^2 = 225$, $\beta^2 = 81$ \textnormal{ και εστίες τα σημεία E($\gamma$,0), E'(-$\gamma$,0). Είναι:}
\[
\gamma^2 = \alpha^2 - \beta^2 = 225 - 81 = 144
\]
\textnormal{Άρα, } $\gamma = 12$. \textnormal{ Επομένως, οι εστίες της έλλειψης είναι:}
\[
\textnormal{E}(12,0), \textnormal{ E}'(-12,0)
\]

\textnormal{β) Εξετάζουμε αν οι συντεταγμένες του σημείου Β επαληθεύουν την εξίσωση της έλλειψης. Για } $x = 0$ \textnormal{ και } $y = 9$ \textnormal{ είναι:}
\[
\frac{0^2}{225} + \frac{9^2}{81} = 1
\]
\textnormal{Η τελευταία ισότητα είναι αληθής, οπότε το σημείο Β είναι σημείο της έλλειψης .}

\textnormal{γ) Η εξίσωση της εφαπτομένης της έλλειψης στο σημείο της B($x_1$,$y_1$) είναι:}
\[
\frac{xx_1}{225} + \frac{yy_1}{81} = 1
\]
\textnormal{Αφού δίνεται ότι B(0,9), η ζητούμενη εξίσωση θα είναι:}
\[
\frac{0x}{225} + \frac{9y}{81} = 1 \textnormal{ ή } y = 9
\]

\textnormal{Η καμπύλη της έλλειψης, οι εστίες E, E' της και η εφαπτομένη ($\varepsilon$) απεικονίζονται στο ακόλουθο σχήμα:}
% TODO: TikZ σχήμα (μην το κατασκευάσεις)
\end{lysh}